\section{Numerical Methods}\label{sec:num_methods}
\subsection{Grid Construction}\label{subsec:grid_construction}

We constructed grids for capital and for productivity.
We will denote said grids as $K$ and $Z$ respectively.
As we had to accommodate the $AR(1)$ process for productivity, we took $Z$ as a baseline grid and formed $K$ to fit the expectations created by $Z$.
Given $Z$, our main benchmark for the $K$ specifications was eliminating (i.e. minimizing) any value maximizing capital at the grid boundaries.

\subsubsection{Z Grid}\label{subsubsec:z_grid}

The $Z$ grid is constructed by the Rouwenhorst method which is a discretization method for approximating continuous state space Markov processes.
The process allow us to create a discrete grid that captures the properties of the underlying continous $AR(1)$ process for productivity.

\subsubsection{K Grid}\label{subsubsec:k_grid}

The $K$ grid's construction depends on $Z$; specifically, we use $k^*(z)$, the optimal capital choice for a given productivity level $z$, to determine the range of our $K$ grid.
We solve the convex cost case $\{F=0, \ \rho_s = 1\}$ and solve for the steady state $k_{ss}$.
The grid is then defined in the interval $[\alpha k_{ss}, \beta k_{ss}]$ where $\alpha$ and $\beta$ parameters are chosen to ensure optimal capital choices never fall outside the grid.
To achieve this, we minimize the point mass of choices at the grid boundaries through the bound scaling parameters:

\begin{equation}
    \begin{aligned}
        \hat{\alpha} &= \underset{\alpha}{\operatorname{argmin}}\ R_{\alpha}:=\frac{\sum \mathbbm{1}_{\{k^*(z) < \alpha k_{ss}\}}}{N} \to \hat{\alpha} = 0.01 \ (R_{\alpha} = 0.00\%)\\
        \hat{\beta} &= \underset{\beta}{\operatorname{argmin}}\ R_{\beta}:=\frac{\sum \mathbbm{1}_{\{k^*(z) > \beta k_{ss}\}}}{N} \to \hat{\beta} = 25 \ (R_{\beta} = 0.20\%)
    \end{aligned}
\end{equation}

\noindent\textbf{NOTE:} We chose not to push $\hat{\beta}$ further than 25 to capture all outliers. We achieved $R_\beta = 0.2\%$ at $\beta = 25$ and tested that the next increment would happen around $\beta=50$. 
At which point we decided accepting $R_\beta = 0.2\%$ as a reasonable tolerance since a marginal reduction in $R_\beta$ would cost increasingly larger $\beta$ specifications.

\vspace{\baselineskip}

With the above result yielding a grid over $[0.01 k_{ss}, 100 k_{ss}]$, we observe multiple orders of magnitude between the grid bounds.
Therefore, we elect to use a non-linear spacing for the $K$ grid to emphasize the lower end while covering the upper extremes.
A $\log$-spaced grid satisfies this requirement and is constructed by defining an auxiliary grid $K_{aux}$ with linear spacing over $[\log(\alpha k_{ss}), \log(\beta k_{ss})]$ and then exponentiating the values to obtain the final $K$ grid:
\begin{equation}
    \begin{aligned}
        K_{aux} &= \{\ln(\alpha k_{ss}), \ln(\alpha k_{ss}) + \Delta, \ln(\alpha k_{ss}) + 2\Delta, \dots, \ln(\beta k_{ss})\} \\
        K &= \{\exp(k_i)\ \forall k_i \in K_{aux}\}
    \end{aligned}
\end{equation}

\subsection{Interpolation}\label{subsec:interpolation}
With a discrete grid, we need to either choose the nearest grid point or interpolate to evaluate the value function at non-grid points.
We choose to use a weighted interpolation of the two nearest grid points for capital.
Interpolation helps sanitize the kinks in value and policy functions, making them more interpretable.

\vspace{\baselineskip}

Our interpolation method normalized the distance from/to the nearest grid points over the total distance between the two grid points encapsulating the optimal value.
Denoting a general case where we want to interpolate a function $f$ at a point $x$ that lies between two grid points $\bar{x}_i$ and $\bar{x}_{i+1}$, the interpolation is given by:

\begin{equation}
    \begin{aligned}
        f(x) &= w_i f(\bar{x}_i) + w_{i+1} f(\bar{x}_{i+1}) \\
        w_i &= \frac{\bar{x}_{i+1} - x}{\bar{x}_{i+1} - \bar{x}_i} \\
        w_{i+1} &= \frac{x - \bar{x}_i}{\bar{x}_{i+1} - \bar{x}_i} \\
    \end{aligned}
\end{equation}


\subsection{Value Function Iteration}\label{subsec:vfi}
The Value Function Iteration (VFI) process is aided by a bellman operator that uses golden-section search to find the optimal capital choice.
Golden-section search refers to an iterative optimization technique for finding the extrema of a unimodal function over some interval $[a, b]$.
The method works by evaluating the function at two interior points, $c$ and $d$, which are determined by the golden ratio $\phi = \frac{1 + \sqrt{5}}{2}$.

\vspace{\baselineskip}

Additionally, we apply a Howard improvement every 5 iterations to speed up convergence.
Howard improvement begins by picking the optimal policy from the previous iteration and evaluating the value function under this policy for $n$ periods (or to infinity by analytical definition).
Afterwards, the policy is updated by finding the optimal choice under the new value function and the process is repeated until convergence.
Howard improvement can significantly reduce the number of iterations in the outer VFI loop necessary for convergence, especially when the policy function stabilizes before the value function does.

\subsection{MoM Estimation Setup}\label{subsec:smm_setup}
We use the given LRD moments as targets and match the moments from a stationary distribution ($\mu$) computed under some parameter composition $\{\gamma_{adj}, F, \rho_s\}$.
The parameters are updated by minimizing the L2 distance between target and output moments using the Nelder-Mead optimization method. 
However, we prefer running the optimization in an unconstrained space, so appropriate transformations are applied to the parameters prior to the optimization.

\begin{equation}
    \begin{aligned}
         \gamma_{adj}\in\ (0, \infty) \to \gamma_{adj}'\in\mathbb{R} &\quad\mid\quad \gamma_{adj}' = \ln(\gamma_{adj}) \\
        F\in\ [0, \infty) \to F'\in\mathbb{R} &\quad\mid\quad F' = \max\{10^{-12},\ \ln(F)\} \\
        \rho_s\in\ [0, 1] \to \rho_s'\in\mathbb{R} &\quad\mid\quad \rho_s' = \operatorname{logit}(\rho_s) = \ln\left(\frac{\rho_s}{1-\rho_s}\right)
    \end{aligned}
\end{equation}

With said transformations, the Nelder-Mead optimization can run unbounded over the set $\{\gamma_{adj}', F', \rho_s'\}$
as long as we use the inverse transformations to compute moments inside the objective function.